\labsection{6}{Build a blockchain from first principles}{sec:lab06-blockchain}

\begin{labproject}{Implement hashing, Merkle roots, chaining, and proof of work --- then break it}
\textbf{Coverage.} Chapter~\ref{ch:blockchain}.\\
\textbf{Goal.} Build the structure Chapter~\ref{ch:blockchain} describes, using a complete standard-library Python program, then detect tampering and explain why a real consensus network requires more than hash links.

\courselabflow{Hash records}{Link blocks}{Mine proof}{Detect tampering}



\end{labproject}

\textbf{Complete worked example.} This runnable fixture uses synthetic inputs. First reproduce it, then adapt the data or ratings with documented evidence. All code below belongs to one file.
\begin{lstlisting}[style=utpterminal]
python all_experiments/lab06_build_a_blockchain.py
\end{lstlisting}

\noindent\textbf{Setup and output folder.}
\lstinputlisting[style=utppython,firstline=1,lastline=10]{all_experiments/lab06_build_a_blockchain.py}

\noindent\textbf{Hash unambiguous serialised data.}
\lstinputlisting[style=utppython,firstline=11,lastline=28]{all_experiments/lab06_build_a_blockchain.py}

\noindent\textbf{Define the block and bounded proof of work.}
\courselabtrace{What one changed transaction affects}{Transaction}{A changed value changes its Merkle root}{Block}{The changed header has a different hash}{Next block}{Its stored previous hash no longer matches}

\lstinputlisting[style=utppython,firstline=29,lastline=47]{all_experiments/lab06_build_a_blockchain.py}

\noindent\textbf{Build the complete chain, including genesis.}
\lstinputlisting[style=utppython,firstline=48,lastline=55]{all_experiments/lab06_build_a_blockchain.py}

\noindent\textbf{Verify every block, then deliberately tamper.}


\textbf{Use the chapter while coding.} Revisit \hyperref[ch:blockchain]{the earlier chapter explanation}; match its operation to the variables and state changes below.

\textbf{Follow the supplied implementation.} The changed transaction alters block 2's Merkle root and hash. Block 3 still stores the old hash. Proof-of-work validity is checked independently; do not rely on a specific extra proof-of-work failure when explaining the guaranteed linkage failure.
\courselabtrace{Trace the exact tampering demonstration}{1. Input / state}{The initial four-block chain passes validate(chain, difficulty) with difficulty 2.}{2. Operation}{Change chain[2]'s transaction to bob->mallory 1000 without updating later stored links.}{3. Check}{Validation must report a broken previous-hash link at block 3; the assertion checks that failure.}


\lstinputlisting[style=utppython,firstline=56,lastline=78]{all_experiments/lab06_build_a_blockchain.py}
\begin{labvisual}{Tampering changes a hash; it does not repair the next link}
\begin{labgraphic}\begin{tikzpicture}[font=\small]
\node[utpstep,text width=2.5cm,minimum height=1.25cm] (a) {Block 0\\previous: zeros\\hash: $H_0$};
\node[utpstep,text width=2.5cm,minimum height=1.25cm,right=.35cm of a] (b) {Block 1\\previous: $H_0$\\hash: $H_1$};
\node[utpstep,fill=UTPPaleOrange,text width=2.5cm,minimum height=1.25cm,right=.35cm of b] (c) {Block 2 changed\\previous: $H_1$\\new hash: $H'_2$};
\node[utpstep,text width=2.5cm,minimum height=1.25cm,right=.35cm of c] (d) {Block 3\\stored previous: $H_2$\\$H_2\ne H'_2$};
\draw[utparrow] (a)--(b);\draw[utparrow] (b)--(c);\draw[utparrow,draw=UTPOrange] (c)--(d);
\end{tikzpicture}\end{labgraphic}\end{labvisual}

\textbf{Check the run.} The following output was captured from the supplied script; it describes the synthetic fixture only.
\lstinputlisting[style=utpterminal]{labs/outputs/lab06.txt}
% Reused from chapters/06-blockchain.tex
\begin{labvisual}{Recall the chapter model}
\begin{labgraphic}
\begin{tikzpicture}[
  c/.style={draw=UTPBlue,fill=UTPPaleBlue,rounded corners=2mm,minimum width=2.2cm,minimum height=0.95cm,align=center,font=\scriptsize\bfseries\color{UTPNavy}},
  auth/.style={draw=UTPNavy,fill=UTPNavy,text=white,rounded corners=2mm,minimum width=2.7cm,minimum height=1.15cm,align=center,font=\scriptsize\bfseries},
  peer/.style={draw=UTPTeal,fill=UTPPaleTeal,circle,minimum size=1.2cm,align=center,font=\scriptsize\bfseries\color{UTPNavy}},
  arr/.style={-{Stealth[length=1.8mm]},thick,draw=UTPMuted}
]
\node[auth] (a) at (-4.3,0) {Central\\authoritative ledger};
\node[c] (u1) at (-4.3,2) {User}; \node[c] (u2) at (-4.3,-2) {User};
\draw[arr](u1)--(a); \draw[arr](u2)--(a);
\node[peer] (p1) at (1.7,1.7) {P1}; \node[peer] (p2) at (4.3,1.7) {P2};
\node[peer] (p3) at (1.7,-1.7) {P3}; \node[peer] (p4) at (4.3,-1.7) {P4};
\draw[arr](p1)--(p2); \draw[arr](p2)--(p4); \draw[arr](p4)--(p3); \draw[arr](p3)--(p1); \draw[arr](p1)--(p4); \draw[arr](p2)--(p3);
\node[font=\scriptsize\color{UTPMuted},align=center] at (3.0,0) {Distributed copies\\+ consensus rules};
\end{tikzpicture}
\end{labgraphic}\end{labvisual}


\begin{verification}
\textbf{Explain and extend your result.} Tamper with one transaction and report the first failed validation check. Explain the downstream work required to rebuild the local chain and why that alone does not model a real consensus attack.

\begin{itemize}\item The script runs from a fresh process and creates the named CSV or report files.\item Every rate uses a stated population and denominator; every model assumption is identified.\item Synthetic outputs are labelled and are not presented as observations of a real institution.\item Change one input, rerun, and explain the direction of the result.\end{itemize}\end{verification}
